\chapter{Звідки все почалося: вибір задачі}\label{ch:g01}

Проєкт почався не з фізики, а з організаційного питання: \emph{яку задачу дати студентам
на хакатоні ФМФ}, щоб за 24--48 годин вони могли зробити щось справжнє, а журі --- чесно
це оцінити. Відповідь, до якої ми прийшли 2026-09-22, --- «реконструкція магнітного
потоку $\psi(R,Z)$ без магнітних датчиків, з метрикою, яка питає, чи це взагалі
рівновага». Цей розділ пояснює, чому саме так, які варіанти ми відкинули і які факти
(E25--E29) та відкриті питання (Q10--Q13) стоять за рішенням. Фізику рівноваги й
величини $\psi$, $\qnf$, $\betaN$, $\li$ див. у \cite{Manual}, розд.~3.

\section{Критерій вибору: де стеля, а де простір}\label{sec:g01-criterion}

Правило, яким ми відсіювали задачі, записано у файлі \texttt{saturated-vs-headroom.md} (тека
\texttt{Our try/01-landscape/}): \emph{задача придатна для хакатону, якщо найкращий
опублікований результат помітно нижчий за очевидну стелю}. Інакше студентська команда за
дві доби нічого не зрушить, а таблиця лідерів перетвориться на лотерею шуму.

\begin{established}{E25 (прочитано)}
Прогноз зриву (disruption) насичений \emph{і} має закриті дані: AUC~0,974 всередині машини;
DisruptionBench --- харнес без даних, архівований 2025-07-21 \cite{RegEst,Spangher2025}.
\end{established}
\noindent Число 0,974 (CCNN на Alcator C-Mod) у наших файлах позначено \texttt{[S]} --- взято зі
зведення пошуку, а не з відкритої статті; факт «архівовано без даних» --- \texttt{[V]}.
Залишковий розрив у прогнозі зриву --- операційний (частота хибних тривог за фіксованого часу
попередження), а не в AUC. Дані C-Mod, DIII-D, EAST не випущені, JET закрито юридично,
KSTAR публічного набору не має (\texttt{02-field-map/E-disruption-why-not.md}). Фізику
зривів див. \cite{Manual}, розд.~10.

Натомість простір знайшовся там, де наявні бенчмарки \emph{не міряють}: фізична
узгодженість (нев'язка рівняння Ґреда--Шафранова), калібрування невизначеності,
обумовленість відображення $\psi\mapsto$ скаляри (розд.~\ref{ch:g02}) і неєдиність
розв'язку (розд.~\ref{ch:g03}). Реконструкція $\psi$ всередині DIII-D \emph{за наявною
метрикою} теж близька до стелі: стартовий набір уже дає $S=0{,}8143$ на 20 відкладених
розрядах, а шлюз Challenge~2 поставлено на 0,85 (\texttt{saturated-vs-headroom.md}, \texttt{[V]}).
Тому ми змінюємо \emph{метрику}, а не задачу.

\section{Рішення D1--D9}\label{sec:g01-decisions}

Реєстр рішень \cite{RegDec} (версія~1, 2026-09-22) має незвичну колонку: «що б його
змінило». Рішення без умови перегляду мовчки успадковується, навіть коли обставини
змінилися.

\begin{center}\small
\begin{tabularx}{\textwidth}{@{}p{0.7cm}Xp{4.3cm}@{}}
\toprule
& Рішення & Що б його змінило \\
\midrule
D1 & ядро --- реконструкція $\psi(R,Z)$, а не прогноз зриву & відкритий датасет зривів із
  мітками $t_{\text{ТГ}}$ \\
D2 & конвеєр рівновага~→ стійкість~→ зрив --- лише в наративі; одна оцінювана задача +
  3--4 номінації на тому самому $\psi$ & горизонт 3+ місяці або 3+ інженери \\
D3 & два треки: базовий (starter kit + таблиця) і дослідницький (журі) & однорідна
  аудиторія \\
D4 & дані sim-to-real: навчання на синтетиці, тест на реальних розрядах & синтетика надто
  «легка» \\
D5 & синтетика з \texttt{FreeGS}/\texttt{FreeGSNKE}/\texttt{TokaLab}, свого генератора не
  пишемо & жоден не дає потрібної геометрії \\
D6 & метрика $S'$ = метрика Sophelio + GS-нев'язка + калібрування UQ & перевірка V-A4
  знайде, що хтось це вже робить \\
D7 & аналіз українською, підсумкові артефакти --- англійською & --- \\
D8 & \textbf{закрито 2026-09-22}: надпровідність не окрема тема; Бін --- опуклий контроль
  у розборі D2 & нічого: питання закрито доказово \\
D9 & розрахунок квенчу магніта --- не оцінювана задача & виміряні дані по котушках під час
  зривів \\
\bottomrule
\end{tabularx}
\end{center}

Два зауваження про історію цих рядків (журнал змін \cite{RegLog}, п.~23 і~26). До
2026-09-29 обґрунтування D1 посилалося на «див.~D9», хоча D9 --- про квенч магніта; тепер
там стоїть посилання на відкинутий варіант «хакатон із прогнозу зриву», E25 і
\texttt{E-disruption-why-not.md}. А англійський підсумок реєстру ще тиждень після
закриття D8 казав, що надпровідність «відкладено до огляду літератури». Обидві неточності
виправлено відкрито, зі старим текстом у журналі.

\paragraph{Відкинуті варіанти.} Реєстр зберігає шість варіантів, від яких ми відмовилися,
разом із причинами:
\begin{enumerate}
\item \emph{Хакатон із прогнозу зриву} (готовий внутрішній протокол) --- насичено, дані
  закриті, а симулятора \texttt{tok-sym-core}, на який спирався протокол, у проєкті немає.
\item \emph{Клон Fusion Equilibrium Challenge} --- змагання живе (Phase~1 до 2026-10-18,
  176~учасників, 684~подання), клонувати надлишково.
\item \emph{Участь у Sophelio командою факультету} --- відхилено: чужий дедлайн відволікає.
\item \emph{Генератор даних на \texttt{fusionsimulator.io}} --- коду локально немає,
  ліцензія суперечлива (Q10), 2,5--3 тижні роботи.
\item \emph{Надпровідність як оцінювана задача} --- немає ground truth; один член втрат
  домінує в $\approx490$ разів, тож «модель» зводиться до підстановки у формулу.
\item \emph{Гіпотеза «Бін~$\leftrightarrow$ вільна межа плазми --- одна математика»} ---
  спростовано власною перевіркою (R2); врятовано переформульовану версію D8
  (розд.~\ref{ch:g03}).
\end{enumerate}
Сюди ж реєстр відносить спростоване твердження «ніхто не нав'язує інтегральну крайову
умову жорстко» (R1, перевірка V-A2; спростовано роботою \cite{McClenaghan2024}). І три визнані ризики, на які план не розраховує:
горизонт менше місяця; формат 24--48~год погано сумісний із дослідницьким треком; команда
--- фактично одна людина плюс 2--4 помічники.

\section{Дані й ліцензії: E26--E29, Q10--Q13}\label{sec:g01-data}

Рішення D4 спирається на відкриті дані, тож першою роботою було \emph{переміряти} все, що
про ці дані пишуть офіційні джерела. Результат --- чотири прочитані факти і чотири зовнішні
питання.

\begin{established}{E26--E29 (прочитано/перевірено)}
\begin{itemize}
\item[E26] Корпус Sophelio --- \textbf{103,86~ГБ, 9\,121 розряд} (7\,041 \texttt{diii\_d\_train}
  + 874 \texttt{diii\_d\_public\_test} + 1\,206 \texttt{mast\_public\_test}), виміряно
  через HF API.
\item[E27] Член Consistency метрики усереднює \textbf{сім} скалярів, не вісім.
\item[E28] Ліцензії: HDB5 --- CC~BY~4.0; ConStellaration --- MIT; Sophelio --- CC~BY~4.0;
  FAIR-MAST --- CC~BY-SA~4.0 із поширенням на похідні бази через §4(b).
\item[E29] Код стартового набору справді MIT; статус \texttt{NOASSERTION} на GitHub
  спричинений дописаним розділом «NOTE ON SCOPE».
\end{itemize}
\end{established}

Чому це не дрібниці. Офіційні числа суперечать самі собі (V-B1, V-B2 \cite{RegVerif}):
на сайті 98~ГБ, у статті 133~ГБ; кількість розрядів в анотації статті й у її тексті
різна. Тому для статистики датасету статтю \cite{FEC} ми вважаємо ненадійною, а
перелічуємо файли самі. Так само з E27: єдине «вісім» --- у застарілій спливній підказці на
сайті, а код скорера (нижче) каже «сім»; для нас код надійніший за сайт.

\begin{center}\small
\begin{tabularx}{\textwidth}{@{}p{0.8cm}Xp{3.3cm}@{}}
\toprule
& Питання & Статус \\
\midrule
Q10 & умови поширення даних, згенерованих \texttt{fusionsimulator.io} (README водночас
  «MIT» і «All rights reserved») & відкрите; нічого похідного не публікуємо \\
Q11 & ліцензія ITPA HDB5 & закрите: CC~BY~4.0 (поле власників порожнє) \\
Q12 & атрибуція скорера й даних Sophelio & закрите: код MIT, дані CC~BY~4.0 \\
Q13 & правова підстава переходу 1\,206 розрядів MAST із CC~BY-SA~4.0 (FAIR-MAST) у випуск
  CC~BY~4.0 & відкрите; частину MAST трактуємо як обтяжену BY-SA \\
\bottomrule
\end{tabularx}
\end{center}

Q13 --- приклад того, як юридичне питання стає науковим обмеженням: переписана версія
гіпотези C8 про провали перенесення на MAST (розд.~\ref{ch:g06}) без MAST-рівноваг із
FAIR-MAST була нездійсненна, а FAIR-MAST несе share-alike. (Того ж дня, 2026-09-29, C8
закрито без нових даних: провали виявилися артефактом обробки знака --- R9.)

\section{Що саме оцінює скорер}\label{sec:g01-scorer}

Метрика конкурсу (версія скорера 3.2.0) --- зважена сума чотирьох членів:
\begin{equation}\label{eq:g01-S}
S = 0{,}55\max(0,R^2_\psi) + 0{,}15\max(0,R^2_{q\beta})
  + 0{,}10\bigl(1-\min(1,D_{\mathrm{LCFS}})\bigr) + 0{,}20\,\mathrm{Consistency},
\end{equation}
де $R^2_\psi$ --- об'єднаний по кадрах коефіцієнт детермінації карти потоку,
$R^2_{q\beta}$ --- середнє $R^2$ для $\qnf$ і $\betaN$, $D_{\mathrm{LCFS}}$ --- відстань
між останніми замкненими поверхнями, а Consistency --- середнє $\max(0,R^2_j)$ по семи
скалярах, \emph{обчислених скорером із поданого $\psi$} тим самим кодом, що й з істини.
Ключова думка дизайнерів: скаляр не можна «підкрутити» окремою головою мережі --- його
заробляє лише $\psi$, яке його справді містить.

\code{fec/starter/fusion_scoring/common.py}{40}{57}{(канал скалярів і ваги метрики)}

\begin{itemize}
\item \emph{Рядки 40--43:} лише два скаляри, $\qnf$ і $\betaN$, подаються окремо --- їх не
  можна відновити з самого $\psi(R,Z)$, бо вони залежать від профілів $F(\psi)$ і
  $p(\psi)$ (\cite{Manual}, розд.~3).
\item \emph{Рядки 45--49:} сім похідних скалярів Consistency --- положення магнітної осі
  $\Rax,\Zax$, витягнутість $\kappa$, верхня й нижня трикутність, об'єм і внутрішня
  індуктивність $\li$. Коментар фіксує, що восьмий, \texttt{dsep}, прибрано у v3: це і є
  E27, звірене з кодом.
\item \emph{Рядки 51--52:} сітка $65\times65$, рядки --- $Z$, стовпці --- $R$.
\item \emph{Рядки 54--57:} ваги \eqref{eq:g01-S}. $R^2_\psi$ домінує (0,55); межа і
  форма в Consistency свідомо перекриваються, тому вага $D_{\mathrm{LCFS}}$ скромна.
\end{itemize}

Саме тут народилася перша гіпотеза проєкту. Дизайнери стверджують, що ідеальна карта
потоку «заробляє все». Але $R^2_\psi$ міряє $\psi$ у нормі $L^2$, а скаляри --- це
положення критичних точок і ліній рівня, тобто \emph{диференціальні} функціонали.
Наскільки \emph{майже} ідеальне $\psi$ дає майже ідеальні скаляри? Це питання
розділу~\ref{ch:g02}.

\section{Концепт: «Магнітна рівновага, якої не видно»}\label{sec:g01-concept}

Концепт v1 (\texttt{Our try/05-concept/\allowbreak hackathon-concept-v1.md}, 2026-09-22) ---
документ для оргкомітету, не протокол для учасників. Коротко:
\begin{itemize}
\item \textbf{Задача.} Вхід: струми котушок полоїдального поля, профілі томсонівського
  розсіяння $T_e$, $n_e$, струм плазми, геометрія. Магнітної діагностики немає. Вихід:
  $\psi$ на сітці $65\times65$, $\qnf$, $\betaN$ і передбачувальні смуги 50/80/90/95\,\%.
  Мотивація --- реакторна: у машинах класу SPARC/ARC нейтрони деградують магнітні датчики
  (питання Q1 до консультантів --- чи це реалістичний сценарій --- відкрите).
\item \textbf{Метрика.}
  $S' = 0{,}40R^2_\psi + 0{,}10R^2_{q\beta} + 0{,}10(1-D_{\mathrm{LCFS}})
  + 0{,}15\,\mathrm{Consistency} + 0{,}15\,\mathrm{GS\_score} + 0{,}10\,\mathrm{Calibration}$.
  Два останні члени --- наш внесок (рішення D6, розд.~\ref{ch:g04}); що $S'$ справді
  змінить упорядкування моделей, для \texttt{GS\_score} перевірено 2026-10-01: стійко не змінює (R12); частина про Calibration --- поки гіпотеза C4.
\item \textbf{Два треки} (D3) і \textbf{номінації}, що споживають те саме $\psi$ (D2):
  фізична узгодженість, чесна невизначеність, робастність, compute-light.
\item \textbf{Дослідницькі підзадачі} взято з наших розборів: спектрально зважена втрата
  (C2; 2026-09-30 перевірено й спростовано~--- R11, розд.~\ref{ch:g02}), L$^2$-регресія на багатозначній цілі (розд.~\ref{ch:g03}),
  одночасні конформні смуги (C6), Deep Ritz проти PINN (C7).
\end{itemize}
Концепт чесно перелічує й ризики: \texttt{Calibration} можна загеймити нескінченно
широкими смугами (член мусить включати гостроту); частину MAST не перевипускаємо (Q13).

\begin{practicum}[перевірка харнесу і умов перегляду]
\begin{enumerate}
\item Переконайтеся, що скорер цілий (E17; потрібен інтернет --- розряди стрімляться з HF): у теці стартового набору виконайте
  \texttt{.venv/bin/python local\_score.py --mode perfect --n-shots 2}, потім те саме з
  \texttt{--mode zeros}. Має вийти рівно $S=1$ і $S=0$. Подивіться, які з семи скалярів
  у режимі \texttt{zeros} мають найгірше $R^2$ --- це підказка до розд.~\ref{ch:g02}.
\item Візьміть будь-яке рішення D1--D9 і сформулюйте \emph{вимірюваний} тест його умови
  перегляду. Наприклад, для D4: яким числом ви б визначили, що синтетика «надто легка»?
\item (Потрібен інтернет.) Відтворіть E26: перелічіть файли датасету через HF API
  (\texttt{HfApi().list\_repo\_tree} з \verb|repo_type="dataset"|, \texttt{recursive=True}, \texttt{expand=True}) і підсумуйте розміри.
  Чи вийде 103,86~ГБ десяткових (96,73~ГіБ)?
\end{enumerate}
\end{practicum}
